\documentclass[12pt]{article}
\usepackage[a4paper,margin=1in]{geometry}\usepackage[T1]{fontenc}
\usepackage{lmodern,microtype}\usepackage{amsmath,amssymb,amsthm,mathtools}
\usepackage{booktabs,array,enumitem,xcolor}\usepackage[numbers,sort&compress]{natbib}
\usepackage[colorlinks=true,linkcolor=blue!55!black,citecolor=blue!55!black,urlcolor=blue!55!black]{hyperref}
\linespread{1.07}
\newtheorem{theorem}{Theorem}[section]\newtheorem{proposition}[theorem]{Proposition}
\newtheorem{lemma}[theorem]{Lemma}\newtheorem{corollary}[theorem]{Corollary}
\theoremstyle{remark}\newtheorem{remark}[theorem]{Remark}
\newcommand{\C}{\mathbb C}\newcommand{\R}{\mathbb R}

\title{The Centered Conjugate-Quartet Shape Manifold\\
\large Exact Coordinates Behind the Physical Quartic Flow}
\author{Liang Wang}\date{July 2026}
\begin{document}\maketitle
\begin{abstract}
Centered-RMS normalization did not contract the physical quartic flow in
RH-212, but it exposed a stronger exact structure.  We classify every
branch-labeled quartet consisting of two conjugate pairs, with zero
barycenter and unit mean square modulus.  In canonical coordinates its roots
are
\[
 \sqrt u\pm i\sqrt{(1-u)(1+\eta)},\qquad
 -\sqrt u\pm i\sqrt{(1-u)(1-\eta)},
\]
where $0\le u\le1$ and $-1\le\eta\le1$.  Its monic polynomial is
\[
 z^4+(2-4u)z^2+4\sqrt u(1-u)\eta z
       +1-\eta^2(1-u)^2.
\]
Consequently the three nontrivial coefficients satisfy the exact algebraic
identity
\[
 c_3^2=4(2-c_2)(1-c_4).
\]

The interior of the branch-labeled quotient is a smooth two-dimensional
surface; explicit inverse formulas recover $(u,\eta)$.  Boundary fibers are
identified: the entire $u=1$ edge collapses to the double pair $\{\pm1\}$,
and branch labels lose orientation at $u=0$.  All 32 physical quartets in the
sixteen-level RH-212 atlas satisfy the manifold identity with maximum
floating residual $8.81\times10^{-15}$; 400 random formula checks have
coefficient error at most $6.67\times10^{-16}$.

This paper provides an exact finite-dimensional quotient, not a small-noise
limit or dynamical law.  Gate A remains open.
\end{abstract}

\section{From a failed contraction to an exact quotient}

RH-212 compared raw, determinant-radius, and centered-RMS coefficient flows
\citep{WangRH212}.  The centered normalization did not reduce adjacent
variation, so it cannot presently be called a renormalization fixed point.
Nevertheless it imposes two exact constraints:
\begin{equation}\label{eq:centerconstraints}
 \sum_{j=1}^{4}q_j=0,
 \qquad \frac14\sum_{j=1}^{4}|q_j|^2=1.
\end{equation}
The inherited real finite operators also force conjugate closure.  These
three facts reduce the apparent three-coefficient family to two real shape
coordinates.

The objective here is purely algebraic: classify that quotient completely,
including its boundary identifications.  No dependence on $\sigma$ is used
in the proof.

\section{Canonical form}

Let a branch-labeled conjugate quartet be
\begin{equation}\label{eq:generalpairs}
 \alpha\pm i b,\qquad \gamma\pm i d,
 \qquad \alpha,\gamma\in\R,\quad b,d\ge0.
\end{equation}
The label records which pair lies on the positive-real and negative-real
branch whenever those branches are distinct.

\begin{lemma}[Centered pair geometry]\label{lem:center}
If the quartet in \eqref{eq:generalpairs} is centered, then
$\gamma=-\alpha$.  After ordering the branches one may write
\begin{equation}\label{eq:abd}
 a\pm ib,\qquad -a\pm id,qquad a,b,d\ge0.
\end{equation}
If it also has unit mean square modulus, then
\begin{equation}\label{eq:rmsrelation}
 a^2+\frac{b^2+d^2}{2}=1.
\end{equation}
\end{lemma}
\begin{proof}
The root sum is $2(\alpha+\gamma)$, proving the first statement.  Relabeling
sets $a=|\alpha|$.  The sum of squared moduli is
$2(a^2+b^2)+2(a^2+d^2)$; division by four gives
\eqref{eq:rmsrelation}.
\end{proof}

\section{Shape coordinates}

Set
\begin{equation}\label{eq:coordinates}
 u=a^2,
 \qquad
 \eta=\frac{b^2-d^2}{b^2+d^2}
\end{equation}
when $b^2+d^2>0$.  Equation \eqref{eq:rmsrelation} gives
\begin{equation}\label{eq:bd}
 b^2=(1-u)(1+\eta),\qquad
 d^2=(1-u)(1-\eta).
\end{equation}
Nonnegativity is equivalent to
\begin{equation}\label{eq:rectangle}
 0\le u\le1,qquad -1\le\eta\le1.
\end{equation}

\begin{theorem}[Canonical shape parameterization]\label{thm:param}
For every $(u,\eta)$ in \eqref{eq:rectangle}, the multiset
\begin{equation}\label{eq:roots}
 \begin{aligned}
 q_{1,2}&=\sqrt u\pm i\sqrt{(1-u)(1+\eta)},\\
 q_{3,4}&=-\sqrt u\pm i\sqrt{(1-u)(1-\eta)}
 \end{aligned}
\end{equation}
is conjugate closed, centered, and has unit mean square modulus.  Conversely,
every quartet of the form \eqref{eq:abd} satisfying
\eqref{eq:rmsrelation} is represented by \eqref{eq:roots}.
\end{theorem}
\begin{proof}
The forward assertions follow by direct substitution.  Conversely set
$u=a^2$ and use \eqref{eq:coordinates}; then
\eqref{eq:rmsrelation} yields \eqref{eq:bd}.
\end{proof}

For $0<u<1$, branch ordering makes the representation unique.  At $u=1$,
$b=d=0$ and every $\eta$ gives the same double pair.  At $u=0$, the two
branches share real part zero and changing $\eta$ to $-\eta$ merely swaps the
two conjugate pairs as an unlabeled multiset.  These are quotient fibers, not
failures of the formula.

\section{Coefficient map and algebraic image}

Let
\begin{equation}
 Q_{u,\eta}(z)=\prod_{j=1}^{4}(z-q_j)
 =z^4+c_2z^2+c_3z+c_4.
\end{equation}

\begin{theorem}[Exact coefficient formulas]\label{thm:coeff}
The coefficient map $\Phi:(u,\eta)\mapsto(c_2,c_3,c_4)$ is
\begin{equation}\label{eq:coefficients}
 \boxed{
 \begin{aligned}
 c_2&=2-4u,\\
 c_3&=4\sqrt u(1-u)\eta,\\
 c_4&=1-\eta^2(1-u)^2.
 \end{aligned}}
\end{equation}
In particular,
\begin{equation}\label{eq:manifoldidentity}
 \boxed{c_3^2=4(2-c_2)(1-c_4).}
\end{equation}
\end{theorem}
\begin{proof}
Factor the two conjugate quadratics:
\begin{equation}
 ((z-a)^2+b^2)((z+a)^2+d^2).
\end{equation}
Expanding and substituting $a^2=u$ and \eqref{eq:bd} gives
\eqref{eq:coefficients}.  Squaring the formula for $c_3$ and using
$2-c_2=4u$ and $1-c_4=\eta^2(1-u)^2$ proves
\eqref{eq:manifoldidentity}.
\end{proof}

The identity alone is not the full semialgebraic description.  Since
$u=(2-c_2)/4$, the image also satisfies
\begin{equation}\label{eq:inequalities}
 -2\le c_2\le2,qquad
 1-\frac{(2+c_2)^2}{16}\le c_4\le1.
\end{equation}
Together with \eqref{eq:manifoldidentity}, these bounds describe the real
coefficient image, with the boundary identifications above.

\section{Inverse formulas and smooth interior}

\begin{proposition}[Interior inverse]\label{prop:inverse}
For $0<u<1$, the branch-labeled coordinates are recovered from the
coefficients by
\begin{equation}\label{eq:inverse}
 u=\frac{2-c_2}{4},
 \qquad
 \eta=\frac{c_3}{4\sqrt u(1-u)}.
\end{equation}
Moreover, $\Phi$ has differential of rank two throughout
$(0,1)\times(-1,1)$.
\end{proposition}
\begin{proof}
The inverse is immediate from \eqref{eq:coefficients}.  The $u$ derivative
has first component $-4$, while the $\eta$ derivative has first component
zero and second component $4\sqrt u(1-u)>0$.  They are linearly independent.
\end{proof}

Thus the normalized physical quartic does not occupy a generic
three-dimensional real coefficient body.  It lies on a smooth
two-dimensional interior surface with explicit compactification.

\section{Finite physical atlas}

The RH-212 atlas contains sixteen scales from $0.04$ to $0.00125$ and two
channels at each scale.  For every endpoint we recover $(u,\eta)$ from the
normalized roots and independently reconstruct the coefficients using
\eqref{eq:coefficients}.

\begin{center}
\begin{tabular}{lr}
\toprule
diagnostic & maximum over 32 endpoints\\
\midrule
manifold identity residual & $8.81\times10^{-15}$\\
coefficient reconstruction error & $8.81\times10^{-15}$\\
coordinate reconstruction error & below $10^{-13}$\\
root multiset reconstruction error & floating roundoff\\
\bottomrule
\end{tabular}
\end{center}

In a separate 400-case random audit over the interior rectangle, direct root
expansion and formula coefficients disagree by at most
$6.67\times10^{-16}$.  These computations verify the implementation; the
theorems themselves are elementary exact algebra.

\section{Interpretation for the divisor-first route}

The result changes the status of centered-RMS normalization.  It is not a
successful contraction in RH-212, but it is the canonical quotient by root
translation and positive scale.  The remaining motion has two distinct
components:
\begin{description}[leftmargin=3.2em,style=nextline]
\item[Axial coordinate $u$]
the squared separation of the two real branch lines;
\item[Asymmetry coordinate $\eta$]
the imbalance between the squared imaginary heights of the positive and
negative real branches.
\end{description}

This vocabulary is intrinsic to the branch-labeled conjugate quartet.  It
does not identify the branches with arithmetic objects and does not produce
an infinite spectrum.

\section{Claim boundary}

Nothing in the algebra proves that $(u_\sigma,\eta_\sigma)$ converges, that
$u_\sigma$ is monotone at all levels, or that the finite points lie on an
autonomous orbit.  It also discards the center and RMS radius, which must be
restored before reconstructing the raw divisor.  Those questions are
separated into later layers.

The exact conclusion is the shape manifold and its coordinate dictionary.
Gate A remains open; Gates B--E, Hilbert--P\'olya, zeta-zero identification,
and RH are untouched.

\bibliographystyle{plainnat}\bibliography{references}\end{document}
